% GENERATED FILE. Edit canonical lessons, then run npm run generate:books.
% canonical-source-hash: fnv1a64:6f544c72e95d7621
% canonical-lessons: ML-C503-paristhiti, ML-C503-parisaram, ML-C503-malinikaranam, ML-C503-malinamakkuka, ML-C503-malinajalam

\chapter{The Environment}
\label{ch:ml-the-environment}
\hlchaptermodality{503}

\begin{chapteropening}
\textbf{By the end of this chapter:} \emph{I can say the environment, surroundings, pollution, to pollute and sewage.}

Complete the last lesson of chapter 503: I can say the environment, surroundings, pollution, to pollute and sewage.
\end{chapteropening}

\section[\texorpdfstring{\ml{പരിസ്ഥിതി} (paristhiti)}{പരിസ്ഥിതി (paristhiti)}]{\ml{പരിസ്ഥിതി} (paristhiti) — the environment}

\label{lesson:ML-C503-paristhiti}



\begin{quote}
Before the new one: say the Malayalam for to fall off, then the Malayalam for to cut down.
\end{quote}

\subsection*{You'll want to know: \ml{പരിസ്ഥിതി}}
\textbf{\ml{പരിസ്ഥിതി}} — \emph{paristhiti} — \textquotedblleft{}the environment\textquotedblright{}.

\begin{quote}
\ml{പ്ലാസ്റ്റിക്} \ml{പരിസ്ഥിതിക്ക്} \ml{വലിയ} \ml{ദോഷമാണ്}.
\end{quote}

\begin{quote}
\emph{plāsṟṟikŭ paristhitikkŭ valiya dōṣamāṇŭ.}
\end{quote}

\begin{quote}
Plastic does great harm to the environment.
\end{quote}

\ml{പരി}-, around, plus \ml{സ്ഥിതി}, a state: the state of things around us.

\begin{cousinweb}
\textbf{\ml{പരിസ്ഥിതി}} (\emph{paristhiti}) is a modern Malayalam word built from Sanskrit \emph{pari}, \textquotedblleft{}around\textquotedblright{}, and \emph{sthiti}, \textquotedblleft{}standing, a state\textquotedblright{} (Monier-Williams).
\end{cousinweb}

\begin{grammarlens}[title={What you've built}]
One more word for getting around.
\end{grammarlens}

\subsection*{Your turn}
\begin{itemize}
\raggedright
  \item \emph{Say it:} \emph{paristhiti}
  \item \emph{Say it:} \emph{paristhiti}, once more
  \item \emph{Say it:} \emph{plāsṟṟikŭ paristhitikkŭ valiya dōṣamāṇŭ.} — Plastic does great harm to the environment.
  \item \emph{Recall:} say the Malayalam for to fall off, then the Malayalam for to cut down, then say \emph{paristhiti} again
\end{itemize}

\begin{tcolorbox}[breakable,colback=teal!4,colframe=teal!35!black,title={Before you move on}]
What does \ml{പരിസ്ഥിതി} mean? (\textquotedblleft{}The environment\textquotedblright{}.) Say it once more.
\end{tcolorbox}
\section[\texorpdfstring{\ml{പരിസരം} (parisaraṁ)}{പരിസരം (parisaraṁ)}]{\ml{പരിസരം} (parisaraṁ) — surroundings; the area around}

\label{lesson:ML-C503-parisaram}



\begin{quote}
Before the new one: say the Malayalam for to cut down, then the Malayalam for the environment.
\end{quote}

\subsection*{You'll want to know: \ml{പരിസരം}}
\textbf{\ml{പരിസരം}} — \emph{parisaraṁ} — \textquotedblleft{}surroundings; the area around\textquotedblright{}.

\begin{quote}
\ml{വീടിന്റെ} \ml{പരിസരം} \ml{വൃത്തിയാക്കണം}.
\end{quote}

\begin{quote}
\emph{vīṭinṟe parisaraṁ vṛttiyākkaṇaṁ.}
\end{quote}

\begin{quote}
The area around the house must be cleaned.
\end{quote}

Notices ask you to keep the \ml{പരിസരം} clean.

\begin{cousinweb}
\textbf{\ml{പരിസരം}} (\emph{parisaraṁ}) is borrowed from Sanskrit \emph{parisara}, \textquotedblleft{}the neighbourhood, the surroundings\textquotedblright{} (Monier-Williams).
\end{cousinweb}

\begin{grammarlens}[title={What you've built}]
One more word for getting around.
\end{grammarlens}

\subsection*{Your turn}
\begin{itemize}
\raggedright
  \item \emph{Say it:} \emph{parisaraṁ}
  \item \emph{Say it:} \emph{parisaraṁ}, once more
  \item \emph{Say it:} \emph{vīṭinṟe parisaraṁ vṛttiyākkaṇaṁ.} — The area around the house must be cleaned.
  \item \emph{Recall:} say the Malayalam for to cut down, then the Malayalam for the environment, then say \emph{parisaraṁ} again
\end{itemize}

\begin{tcolorbox}[breakable,colback=teal!4,colframe=teal!35!black,title={Before you move on}]
What does \ml{പരിസരം} mean? (\textquotedblleft{}Surroundings; the area around\textquotedblright{}.) Say it once more.
\end{tcolorbox}
\section[\texorpdfstring{\ml{മലിനീകരണം} (malinīkaraṇaṁ)}{മലിനീകരണം (malinīkaraṇaṁ)}]{\ml{മലിനീകരണം} (malinīkaraṇaṁ) — pollution}

\label{lesson:ML-C503-malinikaranam}



\begin{quote}
Before the new one: say the Malayalam for the environment, then the Malayalam for surroundings.
\end{quote}

\subsection*{You'll want to know: \ml{മലിനീകരണം}}
\textbf{\ml{മലിനീകരണം}} — \emph{malinīkaraṇaṁ} — \textquotedblleft{}pollution\textquotedblright{}.

\begin{quote}
\ml{നഗരത്തിൽ} \ml{വായു} \ml{മലിനീകരണം} \ml{കൂടുതലാണ്}.
\end{quote}

\begin{quote}
\emph{nagarattil vāyu malinīkaraṇaṁ kūṭutalāṇŭ.}
\end{quote}

\begin{quote}
Air pollution is high in the city.
\end{quote}

\ml{മലിനം}, dirty, plus -\ml{ീകരണം}, making: the making dirty of air, water or land.

\begin{cousinweb}
\textbf{\ml{മലിനീകരണം}} (\emph{malinīkaraṇaṁ}) is built from Sanskrit \emph{malina}, \textquotedblleft{}dirty, impure\textquotedblright{} (Monier-Williams), with \emph{-īkaraṇa}, making into.
\end{cousinweb}

\begin{grammarlens}[title={What you've built}]
One more word for getting around.
\end{grammarlens}

\subsection*{Your turn}
\begin{itemize}
\raggedright
  \item \emph{Say it:} \emph{malinīkaraṇaṁ}
  \item \emph{Say it:} \emph{malinīkaraṇaṁ}, once more
  \item \emph{Say it:} \emph{nagarattil vāyu malinīkaraṇaṁ kūṭutalāṇŭ.} — Air pollution is high in the city.
  \item \emph{Recall:} say the Malayalam for the environment, then the Malayalam for surroundings, then say \emph{malinīkaraṇaṁ} again
\end{itemize}

\begin{tcolorbox}[breakable,colback=teal!4,colframe=teal!35!black,title={Before you move on}]
What does \ml{മലിനീകരണം} mean? (\textquotedblleft{}Pollution\textquotedblright{}.) Say it once more.
\end{tcolorbox}
\section[\texorpdfstring{\ml{മലിനമാക്കുക} (malinamākkuka)}{മലിനമാക്കുക (malinamākkuka)}]{\ml{മലിനമാക്കുക} (malinamākkuka) — to pollute}

\label{lesson:ML-C503-malinamakkuka}



\begin{quote}
Before the new one: say the Malayalam for surroundings, then the Malayalam for pollution.
\end{quote}

\subsection*{You'll want to know: \ml{മലിനമാക്കുക}}
\textbf{\ml{മലിനമാക്കുക}} — \emph{malinamākkuka} — \textquotedblleft{}to pollute\textquotedblright{}.

\begin{quote}
\ml{ഫാക്ടറി} \ml{പുഴ} \ml{മലിനമാക്കി}.
\end{quote}

\begin{quote}
\emph{phākṭaṟi puḻa malinamākki.}
\end{quote}

\begin{quote}
The factory polluted the river.
\end{quote}

Past: \ml{മലിനമാക്കി}, \emph{malinamākki}.

\ml{മലിനം}, dirty, plus \ml{ആക്കുക}, to make. The everyday word is \ml{വൃത്തികേടാക്കുക}.

\begin{cousinweb}
\textbf{\ml{മലിനമാക്കുക}} (\emph{malinamākkuka}) is a Malayalam verb built from \ml{മലിനം}, dirty, borrowed from Sanskrit \emph{malina} (Monier-Williams), and \ml{ആക്കുക}, to make.
\end{cousinweb}

\begin{grammarlens}[title={What you've built}]
One more word for getting around.
\end{grammarlens}

\subsection*{Your turn}
\begin{itemize}
\raggedright
  \item \emph{Say it:} \emph{malinamākkuka}
  \item \emph{Say it:} \emph{malinamākkuka}, once more
  \item \emph{Say it:} \emph{phākṭaṟi puḻa malinamākki.} — The factory polluted the river.
  \item \emph{Recall:} say the Malayalam for surroundings, then the Malayalam for pollution, then say \emph{malinamākkuka} again
\end{itemize}

\begin{tcolorbox}[breakable,colback=teal!4,colframe=teal!35!black,title={Before you move on}]
What does \ml{മലിനമാക്കുക} mean? (\textquotedblleft{}To pollute\textquotedblright{}.) Say it once more.
\end{tcolorbox}
\section[\texorpdfstring{\ml{മലിനജലം} (malinajalaṁ)}{മലിനജലം (malinajalaṁ)}]{\ml{മലിനജലം} (malinajalaṁ) — sewage; dirty water}

\label{lesson:ML-C503-malinajalam}



\begin{quote}
Before the new one: say the Malayalam for pollution, then the Malayalam for to pollute.
\end{quote}

\subsection*{You'll want to know: \ml{മലിനജലം}}
\textbf{\ml{മലിനജലം}} — \emph{malinajalaṁ} — \textquotedblleft{}sewage; dirty water\textquotedblright{}.

\begin{quote}
\ml{മലിനജലം} \ml{നേരിട്ട്} \ml{കായലിലേക്ക്} \ml{ഒഴുകുന്നു}.
\end{quote}

\begin{quote}
\emph{malinajalaṁ nēriṭṭŭ kāyalilēkkŭ oḻukunnu.}
\end{quote}

\begin{quote}
Sewage flows straight into the backwaters.
\end{quote}

\ml{മലിനം}, dirty, plus \ml{ജലം}, water, the formal word for \ml{വെള്ളം}.

\begin{cousinweb}
\textbf{\ml{മലിനജലം}} (\emph{malinajalaṁ}) is a Malayalam compound built from two Sanskrit words, \emph{malina}, \textquotedblleft{}dirty\textquotedblright{}, and \emph{jala}, \textquotedblleft{}water\textquotedblright{} (Monier-Williams).
\end{cousinweb}

\begin{grammarlens}[title={What you've built}]
One more word for getting around.
\end{grammarlens}

\subsection*{Your turn}
\begin{itemize}
\raggedright
  \item \emph{Say it:} \emph{malinajalaṁ}
  \item \emph{Say it:} \emph{malinajalaṁ}, once more
  \item \emph{Say it:} \emph{malinajalaṁ nēriṭṭŭ kāyalilēkkŭ oḻukunnu.} — Sewage flows straight into the backwaters.
  \item \emph{Recall:} say the Malayalam for pollution, then the Malayalam for to pollute, then say \emph{malinajalaṁ} again
\end{itemize}

\begin{tcolorbox}[breakable,colback=teal!4,colframe=teal!35!black,title={Before you move on}]
What does \ml{മലിനജലം} mean? (\textquotedblleft{}Sewage; dirty water\textquotedblright{}.) Say it once more.
\end{tcolorbox}
